\subsection{Classical Planning}
\label{sec:foundations_planning}

Russell and Norvig~\cite{Russell2016} describe classical planning as operating on factored representations, meaning the world state is decomposed into individual variables like ``robot at location A'' or ``door is open'' rather than treating each complete world situation as an indivisible atomic state.

\begin{defn}[Planning Problem~\cite{Bylander1994}]
    \label{defn:planning_problem}
    A planning problem is a 4-tuple $\langle P, O, I, G \rangle$ where $P$ is a finite set of atomic propositions, $O$ is a finite set of operators, $I \subseteq P$ is the initial state, and $G \subseteq P$ is the goal specification.
\end{defn}

\begin{defn}[State~\cite{Russell2016}]
    \label{defn:state}
    A state $s$ is a subset $s \subseteq P$ representing the propositions that are true in that state.
\end{defn}

States follow database semantics with two key assumptions~\cite{Russell2016,Bylander1994}. The \emph{closed-world assumption} treats unmentioned propositions as false: any proposition $p \in P$ not explicitly in $s$ is considered false in state $s$, eliminating the need to explicitly represent negative facts. The \emph{unique names assumption} ensures that distinct object names refer to distinct entities in the domain.

\begin{defn}[Operator~\cite{Bylander1994}]
    \label{defn:operator}
    An operator $o \in O$ is a triple $\langle \text{precond}(o), \text{add}(o), \text{delete}(o) \rangle$ where $\text{precond}(o) \subseteq P$ are the preconditions, $\text{add}(o) \subseteq P$ are the add effects, and $\text{delete}(o) \subseteq P$ are the delete effects.
\end{defn}

Operators are ground, meaning all variables have been replaced by specific objects~\cite{Bylander1994}.

\begin{defn}[Applicability and Result~\cite{Bylander1994}]
    \label{defn:applicability}
    An operator $o$ is applicable in state $s$ if and only if $\text{precond}(o) \subseteq s$. The result of applying $o$ in $s$ is:
    \begin{equation*}
        \text{RESULT}(s,o) = (s \setminus \text{delete}(o)) \cup \text{add}(o)
    \end{equation*}
\end{defn}

\begin{example}[Operator Application]
    \label{ex:operator_application}
    Consider state $s = \{a, b\}$ and operator $o = \langle \{a\}, \{c\}, \{b\} \rangle$ with $\text{precond}(o) = \{a\}$, $\text{add}(o) = \{c\}$, and $\text{delete}(o) = \{b\}$. Since $\{a\} \subseteq \{a, b\}$, operator $o$ is applicable in $s$. The result is $\text{RESULT}(s, o) = (\{a, b\} \setminus \{b\}) \cup \{c\} = \{a, c\}$.
\end{example}

\begin{defn}[Plan~\cite{Russell2016}]
    \label{defn:plan}
    A plan is a sequence of operators $\langle o_1, o_2, \ldots, o_k \rangle$ such that each $o_i$ is applicable in state $s_{i-1}$, where $s_0 = I$ and $s_i = \text{RESULT}(s_{i-1}, o_i)$, and the final state satisfies $G \subseteq s_k$.
\end{defn}

\begin{defn}[Solvability~\cite{Bylander1994}]
    \label{defn:solvability}
    A planning problem $\langle P, O, I, G \rangle$ is solvable if there exists a plan; it is unsolvable otherwise.
\end{defn}

Determining whether a planning problem is solvable is PSPACE-complete~\cite{Bylander1994}, which has motivated extensive research into efficient algorithms and search techniques.

While the preceding definitions establish the formal semantics of planning, practical planning systems require a concrete specification language. The \ac{PDDL} was developed to provide a standardized language for describing planning problems, initially created for the 1998 \ac{ICAPS} competition~\cite{Ghallab1998}. \ac{PDDL} has since become the de facto standard for representing classical planning problems, evolving through multiple versions to support increasingly expressive features while maintaining a foundational core based on the \ac{STRIPS} representation~\cite{Fikes1971}. The \ac{STRIPS} framework established the foundational principles for modeling planning domains where actions have deterministic effects and the world is fully observable, which \ac{PDDL} extends with additional expressiveness and standardization.

Throughout this thesis, the term ``propositional \ac{PDDL}'' refers to \ac{PDDL} specifications using the \texttt{:strips} and \texttt{:typing} requirements without numeric fluents or other extensions from \ac{PDDL} 2.1 and later versions. The \texttt{:strips} requirement indicates the use of basic precondition-effect actions following the classical \ac{STRIPS} model, while \texttt{:typing} enables hierarchical type declarations for objects, allowing parameters to be restricted to specific types. Another common requirement is \texttt{:equality}, which enables the use of equality predicates to compare object references. This thesis focuses on \texttt{:strips} and \texttt{:typing}, corresponding to the classical planning model with propositional state representations, typed objects, and deterministic actions with conjunctive preconditions and effects.

\ac{PDDL} separates planning specifications into two complementary files: a domain file that defines the reusable structure of a planning domain, and a problem file that specifies a particular instance within that domain. This separation enables domain definitions to be reused across multiple problem instances, similar to the distinction between schema and data in database systems.

A domain file specifies the abstract structure using the following components:
\begin{equation*}
    \text{Domain} = \langle \text{Requirements}, \text{Predicates}, \text{Actions} \rangle
\end{equation*}
where $\text{Requirements}$ declares which \ac{PDDL} features are used (such as \texttt{:strips}, \texttt{:typing}, \texttt{:equality}), $\text{Predicates}$ defines the vocabulary of Boolean properties, and $\text{Actions}$ specifies the available operators.

A problem file instantiates the domain with concrete elements:
\begin{equation*}
    \text{Problem} = \langle \text{DomainRef}, \text{Objects}, \text{Init}, \text{Goal} \rangle
\end{equation*}
where $\text{DomainRef}$ references the domain name, $\text{Objects}$ enumerates the entities in this instance, $\text{Init}$ specifies the initial state, and $\text{Goal}$ defines the goal specification. For instance, a logistics domain file would define predicates like $\text{at}(x, y)$ and actions like $\text{drive}$, while a problem file would specify particular trucks, locations, and delivery goals.

Predicates define the vocabulary of a planning domain as parameterized Boolean functions. Formally, a predicate declaration has the form:
\begin{equation*}
    \text{pred}(x_1, x_2, \ldots, x_n)
\end{equation*}
where $\text{pred}$ is the predicate name and $x_1, \ldots, x_n$ are typed or untyped parameters. When instantiated with concrete objects, a predicate becomes a ground atom (proposition) that can be true or false in a state.

For example, in the blocksworld domain\footnote{Blocksworld is a classical planning benchmark involving stacking and unstacking blocks on a table; see Example~\ref{ex:blocksworld} for a detailed illustration.}, the predicate $\text{on}(x, y)$ with parameters $x$ and $y$ representing blocks becomes the ground atom $\text{on}(a, b)$ when instantiated with specific blocks $a$ and $b$, representing ``block $a$ is on block $b$.'' The set of all ground atoms formed from domain predicates and problem objects constitutes the proposition set $P$ in the planning problem formalization.

\ac{PDDL} actions define parameterized operators using a precondition-effect structure. An action schema has the form:
\begin{equation*}
    \text{action}(\text{params}) = \langle \text{precond}(\text{params}), \text{eff}(\text{params}) \rangle
\end{equation*}
where $\text{params}$ are typed parameters, $\text{precond}(\text{params})$ is a logical formula specifying preconditions, and $\text{eff}(\text{params})$ specifies effects as a conjunction of positive and negative literals.

In basic \ac{PDDL}, preconditions are conjunctions of positive atoms, while effects are conjunctions of positive and negated atoms. The set-based operator representation directly captures this structure:
\begin{align*}
    \text{precond}(o) & = \{p \mid p \text{ appears positively in the precondition of } o\} \\
    \text{add}(o)     & = \{p \mid p \text{ appears positively in the effect of } o\}       \\
    \text{delete}(o)  & = \{p \mid \neg p \text{ appears in the effect of } o\}
\end{align*}

Ground operators result from instantiating action parameters with specific objects from the problem definition, yielding the finite operator set $O$ in the planning problem tuple.

The initial state in \ac{PDDL} is specified as a conjunction of ground atoms:
\begin{equation*}
    \text{Init} = \bigwedge_{i=1}^{n} p_i
\end{equation*}
which corresponds to the set $I = \{p_1, p_2, \ldots, p_n\} \subseteq P$ under the closed-world assumption where unmentioned propositions are false.

Goal specifications in \ac{PDDL} can be arbitrary logical formulas in the full language, but in basic \ac{PDDL} with \texttt{:strips} requirements they are restricted to conjunctions of positive atoms. Formally, a \ac{PDDL} goal $\text{Goal} = g_1 \land g_2 \land \cdots \land g_m$ corresponds to the goal set $G = \{g_1, g_2, \ldots, g_m\} \subseteq P$ in the planning problem formalization. A state $s$ satisfies the goal if and only if $G \subseteq s$. Note that negative goal literals (e.g., $\neg p$) require \ac{PDDL} 3.0 extensions~\cite{Gerevini2005}; this thesis focuses on positive goal conjunctions.

While \ac{PDDL} has evolved through multiple versions, \ac{PDDL} 2.1~\cite{Fox2003} added numeric fluents and durative actions, \ac{PDDL} 2.2~\cite{Edelkamp2004} introduced derived predicates and timed initial literals, \ac{PDDL} 3.0~\cite{Gerevini2005} added preferences and constraints, and \ac{PDDL}+~\cite{Fox2002} supports continuous processes. This thesis focuses on basic forms of \ac{PDDL} due to their foundational simplicity and well-understood semantics. Simple \ac{PDDL} formulations with conjunctive preconditions and effects provide sufficient expressiveness for systematic adaptation of inconsistency measures without the complications of temporal reasoning, complex preferences, or soft constraints.

\begin{example}[Blocksworld Domain]
    \label{ex:blocksworld}
    Consider a simplified blocksworld domain with predicates $\text{on}(x,y)$, $\text{ontable}(x)$, $\text{clear}(x)$, and $\text{handempty}$. The action for stacking block $x$ on block $y$ is defined as:
    \begin{align*}
         & \text{action } \text{stack}(x, y)                                             \\
         & \quad \text{precond}: \text{holding}(x) \land \text{clear}(y)                 \\
         & \quad \text{eff}: \text{on}(x,y) \land \text{clear}(x) \land \text{handempty} \\
         & \qquad\qquad \land \neg\text{holding}(x) \land \neg\text{clear}(y)
    \end{align*}

    When instantiated with specific blocks $a$ and $b$ from a problem file, this yields the ground operator:
    \begin{align*}
        \text{stack}(a,b) = \langle & \{\text{holding}(a), \text{clear}(b)\},                \\
                                    & \{\text{on}(a,b), \text{clear}(a), \text{handempty}\}, \\
                                    & \{\text{holding}(a), \text{clear}(b)\} \rangle
    \end{align*}
    demonstrating how \ac{PDDL} action schemas correspond to the set-based operator representation.
\end{example}

This formalization of \ac{PDDL} provides the necessary foundation for discussing how planning problems are specified in practice, while the earlier formal planning definitions establish the semantic interpretation of these specifications as state spaces and reachability problems.

Modern planning systems employ various search strategies to explore different ways of reaching goals, such as searching forward from the initial state, working backward from the goal, or examining the space of possible plans~\cite{Ghallab2004,Russell2016}. These strategies include techniques that build graphical representations of planning problems to identify conflicts and estimate solution difficulty~\cite{Blum1997}, and heuristic search methods~\cite{Bonet2001,Hoffmann2001} that use simplified versions of problems to guide their search toward solutions. These heuristic approaches have significantly improved planning performance for solvable instances, but when these systems encounter an unsolvable problem, they typically provide minimal diagnostic information beyond reporting failure~\cite{Eriksson2018}. This diagnostic gap represents a significant limitation for users attempting to understand the source of conflicts or determine appropriate modifications to achieve solvability.
